% =============================================================
% 08_crossdomain.tex — 分野横断の整理と試験戦略
% =============================================================
\chapter{分野横断の整理と試験戦略}

\chaptermeta{%
  \item 横断テーマ（法律/セキュリティ/マネジメント/AI）を整理できる
  \item ひっかけ選択肢のパターンを見抜ける
  \item 本試験の時間配分を設計できる
  \item 捨て問の判断基準を持てる
}{45分}{3}{第1--7章すべて}

本章は知識のインプットではなく「戦略」の章です。
7章までに学んだ内容を\textbf{横断的に整理}し、
本試験で確実に合格点を取る\textbf{戦術}を身につけます。

% =============================================================
\section{横断テーマ4つ}
% =============================================================

\begin{teigi}[分野をまたぐ4大テーマ]
Iパスの出題は3分野に分かれますが、次の4テーマは分野を横断して出題されます。
\begin{enumerate}
  \item \termtag{法律・制度}：個人情報保護法、著作権法、不正アクセス禁止法、労働者派遣法
  \item \termtag{セキュリティ}：CIA三要素、暗号、認証（ストラテジ系でも出る）
  \item \termtag{マネジメント手法}：PDCAサイクル、リスク管理（全分野共通）
  \item \termtag{AI/データ活用}：AI、ビッグデータ、IoT（全分野から出題）
\end{enumerate}
\end{teigi}

\begin{hikaku}[横断テーマ：どの章で学んだか]
\comptable{テーマ & 主に学んだ章 & 他の章での出題例}{
法律・制度 & 第2章（法務） & 第7章（不正アクセス禁止法） \\
セキュリティ & 第7章 & 第2章（情報倫理）、第4章（監査） \\
マネジメント手法 & 第3--4章 & 第1章（経営戦略としてのPDCA） \\
AI/データ活用 & 第7章（AI概論） & 第1章（DX）、第5章（アルゴリズム）
}
\end{hikaku}

\subsection{法律の横断整理}

\begin{hikaku}[法律の出題マップ]
\comptable{法律 & 保護対象/規制内容 & 関連する章}{
個人情報保護法 & 個人情報の適正な取扱い & 第2章・第7章 \\
著作権法 & 創作的な表現 & 第2章 \\
特許法 & 技術的発明 & 第2章 \\
不正アクセス禁止法 & ID/PWの不正使用 & 第2章・第7章 \\
労働者派遣法 & 派遣労働者の保護 & 第2章 \\
特定電子メール法 & 広告メールのオプトイン & 第2章
}
\end{hikaku}

% =============================================================
\section{ひっかけ選択肢の見抜き方}
% =============================================================

\begin{teigi}[ひっかけパターン5選]
\begin{enumerate}
  \item \textbf{似た用語の入れ替え}：SLAとSLM、DNSとDHCP、TCPとUDP——見た目が似た用語を入れ替えて不正解肢に
  \item \textbf{主語のすり替え}：「監査人が改善を実施する」→ 実施は被監査部門。主語に注目
  \item \textbf{「すべて」「必ず」の断定}：「OSSはすべて無料」「暗号化すれば必ず安全」→ 例外がある表現は疑う
  \item \textbf{正しい説明＋別の用語}：「短期間で反復開発→ウォーターフォール」→ 説明は正しいが用語が違う
  \item \textbf{部分的に正しい選択肢}：前半は正しく後半が誤り → 最後まで読むこと
\end{enumerate}
\end{teigi}

\begin{pitfall}[最初の選択肢で安心しない]
\begin{itemize}
  \item 選択肢アが正しそうに見えても、必ずイ〜エまで確認する
  \item 「より適切なもの」を聞かれている場合、アも正しいがイがもっと正確、ということがある
  \item 消去法が有効——明らかに誤りの選択肢を先に消す
\end{itemize}
\end{pitfall}

% =============================================================
\section{本試験の時間配分}
% =============================================================

\begin{teigi}[120分100問の時間配分]
\begin{itemize}
  \item 1問あたり約\textbf{72秒}（120分 ÷ 100問）
  \item \textbf{ストラテジ系}（問1〜35）：約40分
  \item \textbf{マネジメント系}（問36〜55）：約25分
  \item \textbf{テクノロジ系}（問56〜100）：約45分
  \item \textbf{見直し}：10分を確保
\end{itemize}
\end{teigi}

\begin{advice}[時間配分のコツ]
\begin{itemize}
  \item 30秒考えてわからない問題は「仮回答」してマークし、後で戻る
  \item 計算問題は時間がかかるので、知識問題を先に片付ける
  \item 最後の10分で「マーク忘れ」「問題番号のずれ」を確認
\end{itemize}
\end{advice}

\subsection{捨て問の判断基準}

\begin{teigi}[捨て問の判断]
\begin{itemize}
  \item \textbf{捨てるべき問題}：2分以上悩んでも解けない計算問題、見たことのない新用語
  \item \textbf{捨ててはいけない問題}：知識系の問題（消去法で絞れる可能性がある）
  \item 合格基準は600/1000点。100問中60問正解が目安
  \item ただし各分野で300点以上が必要 → 得意分野だけでは合格できない
\end{itemize}
\end{teigi}

% =============================================================
\section{直前チェック10パターン}
% =============================================================

\begin{matome}[試験直前に確認する10パターン]
\begin{enumerate}
  \item \textbf{PLの5段階利益}：売上総→営業→経常→税引前→当期純
  \item \textbf{知的財産権の保護期間}：特許20年、実用新案10年、意匠25年、商標10年（更新可）、著作権死後70年
  \item \textbf{請負/準委任/派遣}：指揮命令権と完成責任の有無
  \item \textbf{V字モデル}：要件定義↔運用テスト、外部設計↔システムテスト
  \item \textbf{MTBF/MTTR/稼働率}：稼働率 $= \text{MTBF} / (\text{MTBF} + \text{MTTR})$
  \item \textbf{TCP vs UDP}：TCP=信頼性、UDP=速度
  \item \textbf{正規化}：第1=繰り返し除去、第2=部分従属除去、第3=推移従属除去
  \item \textbf{公開鍵暗号 vs デジタル署名}：暗号は受信者の公開鍵、署名は送信者の秘密鍵
  \item \textbf{CIA三要素}：機密性・完全性・可用性
  \item \textbf{AI/ML/DL}：AI $\supset$ ML $\supset$ DL の入れ子
\end{enumerate}
\end{matome}

% =============================================================
\section{過去問ドリル}
% =============================================================

\shoumatsuhead{過去問ドリル：横断問題}

\begin{itpmon}[\sourcebadge{R7}\enspace\diffbadge{標}]{%
A社がシステム開発をB社に委託し、B社がC社に再委託している。C社の作業者に対してA社の社員が直接作業指示を行うことは、どのような観点から問題があるか。}
\sentakushi
  {著作権法に違反する可能性がある}
  {偽装請負にあたり、労働者派遣法に違反する可能性がある}
  {個人情報保護法に違反する可能性がある}
  {独占禁止法に違反する可能性がある}
\end{itpmon}
\begin{kaisetsu}{イ}
請負・準委任契約で発注者が直接指揮命令すると\textbf{偽装請負}。法律的には\textbf{労働者派遣法}違反。再委託先への直接指示も同様に問題。第2章の契約形態と第8章の横断視点がつながる。
\end{kaisetsu}

\begin{itpmon}[\sourcebadge{original}\enspace\diffbadge{強}]{%
次の選択肢のうち、記述の内容と用語の組合せが正しいものはどれか。}
\sentakushi
  {データを暗号化して身代金を要求するマルウェア — フィッシング}
  {1〜4週間の短い反復で開発する手法 — ウォーターフォール}
  {サービス品質を数値で取り決めた合意文書 — SLA}
  {故障時に機能を縮退して運転を継続する設計 — フェールセーフ}
\end{itpmon}
\begin{kaisetsu}{ウ}
アはランサムウェアの説明にフィッシングの用語を当てた誤り。イはアジャイルの説明にウォーターフォールを当てた誤り。ウは正しい（SLA=サービスレベル合意）。エはフェールソフトの説明にフェールセーフを当てた誤り。典型的な「説明と用語の入れ替え」ひっかけ。
\end{kaisetsu}

\begin{itpmon}[\sourcebadge{original}\enspace\diffbadge{標}]{%
本試験で残り30分、未回答が20問ある場合の対処として、最も適切なものはどれか。}
\sentakushi
  {すべて空欄のまま見直しに時間を使う}
  {未回答の問題を順に解き、わからない問題も仮回答してマークする}
  {最初の問題から解き直す}
  {難しそうな問題を先に解き、簡単な問題は後回しにする}
\end{itpmon}
\begin{kaisetsu}{イ}
空欄は0点確定。\textbf{仮回答でもマーク}すれば25\%の確率で正解。残り時間が少ないときは「全問マーク」を最優先し、余った時間で見直す。
\end{kaisetsu}

% =============================================================
\section{よくあるミスと到達確認}
% =============================================================

\begin{yokumiss}[この章でよくあるミス]
\begin{enumerate}[leftmargin=1.6em, itemsep=5pt]
  \item \textbf{横断テーマを見落とす}——セキュリティはストラテジ系でも出題される。
  \item \textbf{ひっかけの「すべて」「必ず」に引っかかる}——断定表現は疑う。
  \item \textbf{時間切れで空欄を残す}——空欄は0点。仮回答でもマークすべき。
  \item \textbf{得意分野だけで合格しようとする}——各分野300点以上の足切りがある。
  \item \textbf{計算問題に時間をかけすぎる}——2分以上悩んだら仮回答して先へ進む。
\end{enumerate}
\end{yokumiss}

\begin{tatsaku}
  \item 横断テーマ4つを言え、それぞれの出題例を挙げられる
  \item ひっかけパターン5選をすべて説明できる
  \item 本試験の分野別時間配分を設計できる
  \item 捨て問の判断基準を持っている
  \item 直前チェック10パターンを全部暗唱できる
  \item 各分野の足切りライン（300点）を理解している
  \item 消去法のテクニックを使える
  \item 仮回答→後で見直し の戦略を実行できる
\end{tatsaku}
